\documentclass[11pt]{article}
\usepackage[margin=1in]{geometry}
\usepackage{amsmath,amssymb}
\usepackage{booktabs}
\usepackage{array}
\usepackage{tabularx}
\usepackage[hidelinks]{hyperref}

\newcommand{\avg}[1]{\langle #1 \rangle}
\newcommand{\neff}{n_e}
\newcommand{\nc}{n_c}
\newcommand{\Sf}{S_f}
\newcommand{\So}{S_0}

\title{Roughness-Scale Project --- Key Equations}
\author{}
\date{}

\begin{document}
\maketitle

%% ======================================================================
\section{Saint-Venant Equations with Manning Friction}
%% ======================================================================

\subsection{Shallow Water Equations (1-D)}

Continuity:
\begin{equation}\label{eq:continuity}
  \frac{\partial h}{\partial t} + \frac{\partial (Uh)}{\partial x} = p - i
\end{equation}

Momentum:
\begin{equation}\label{eq:momentum}
  \frac{\partial U}{\partial t} + U\frac{\partial U}{\partial x}
  + g\frac{\partial h}{\partial x} + g\bigl(\Sf - \So\bigr) = 0
\end{equation}

\subsection{Manning's Friction Closure}

Friction slope:
\begin{equation}\label{eq:Sf}
  \Sf = \frac{n^2\, U^2}{h^{4/3}}
\end{equation}

Unit discharge (kinematic wave):
\begin{equation}\label{eq:q}
  q = \frac{1}{n}\, h^{5/3}\, \So^{1/2}
\end{equation}


%% ======================================================================
\section{Equivalent Roughness: the Effect Ratio}
%% ======================================================================

The \emph{effect ratio} $\neff/\avg{n}$ is the equivalent Manning's~$n$ normalised by the spatial mean.
A value below unity means the heterogeneous surface behaves as if it were smoother than the average roughness.

\subsection{Lotter Equation (Equal-Discharge Condition)}

The equal-discharge condition requires that total discharge from the heterogeneous surface match that from a uniform surface with $\neff$:
\begin{equation}\label{eq:lotter}
  \neff = \frac{\avg{h}^{5/3}}{\avg{\dfrac{h^{5/3}}{n}}}
\end{equation}

Dividing by $\avg{n}$ gives the effect ratio in Lotter form:
\begin{equation}\label{eq:lotter-ratio}
  \frac{\neff}{\avg{n}}
  = \frac{\avg{h}^{5/3}}{\avg{n}\;\avg{\dfrac{h^{5/3}}{n}}}
\end{equation}

\subsection{Jensen's Inequality}

Because $h^{5/3}$ is convex:
\begin{equation}\label{eq:jensen}
  \avg{h^{5/3}} > \avg{h}^{5/3}
\end{equation}
This is the fundamental reason $\neff/\avg{n} < 1$: vegetation induces depth variance; the nonlinear weighting $h^{5/3}$ means fast (deep, smooth) zones carry disproportionately more discharge.


%% ======================================================================
\section{Spatial Decomposition of $\Sf$}
%% ======================================================================

\subsection{Reynolds Decomposition}

Decompose the spatially varying fields into mean and fluctuation:
\begin{equation}
  n = \avg{n} + n', \quad
  U = \avg{U} + U', \quad
  h = \avg{h} + h'
\end{equation}

Normalised fluctuations:
\begin{equation}
  r = \frac{n'}{\avg{n}}, \qquad
  \upsilon = \frac{U'}{\avg{U}}, \qquad
  \eta = \frac{h'}{\avg{h}}
\end{equation}

\subsection{Second-Order Expansion of Spatially Averaged $\Sf$}

\begin{equation}\label{eq:Sf-decomp}
  \frac{\avg{\Sf}\;\avg{h}^{4/3}}{\avg{n}^2\;\avg{U}^2}
  \;\approx\;
  1
  + \avg{r^2}
  + \avg{\upsilon^2}
  + 4\avg{r\upsilon}
  - \tfrac{8}{3}\bigl(\avg{r\eta} + \avg{\upsilon\eta}\bigr)
  + \tfrac{14}{9}\avg{\eta^2}
\end{equation}

\begin{center}
\begin{tabular}{@{}ll@{}}
\toprule
Term & Interpretation \\
\midrule
$\avg{r^2}$            & Roughness variance \\
$\avg{\upsilon^2}$      & Velocity variance (fast/slow zones) \\
$4\avg{r\upsilon}$      & Roughness--velocity covariance \\
$-\tfrac{8}{3}\avg{r\eta}$ & Depth--roughness covariance \\
$-\tfrac{8}{3}\avg{\upsilon\eta}$ & Depth--velocity covariance \\
$\tfrac{14}{9}\avg{\eta^2}$ & Depth variance \\
\bottomrule
\end{tabular}
\end{center}


%% ======================================================================
\section{Correction Factor: Second-Order Taylor Expansion}
%% ======================================================================

The effect ratio expanded to second order in $(r, \upsilon, \eta)$:
\begin{equation}\label{eq:CF1}
  \frac{\neff}{\avg{n}}
  \;\approx\;
  1
  - \avg{r^2}
  - \tfrac{5}{9}\avg{\eta^2}
  + \tfrac{5}{3}\avg{r\eta}
\end{equation}

Alternative (velocity-based):
\begin{equation}\label{eq:CF2}
  \frac{\neff}{\avg{n}}
  \;\approx\;
  1
  + \tfrac{1}{4}\avg{r^2}
  + \tfrac{5}{4}\avg{\upsilon^2}
  + \tfrac{5}{2}\avg{r\upsilon}
\end{equation}

\subsection{Prediction Hierarchy}

\begin{description}
  \item[$T_0$] Leading term: uses only roughness variance $\avg{r^2}$ (known from the $n$-field alone).
  \item[$T_1$] Adds roughness--velocity and roughness--depth covariances.
  \item[$T_2$] Full 2nd-order: all six terms --- best accuracy but requires simulation fields.
  \item[CF approx] Replaces simulation covariances with analytical correction factors.
\end{description}


%% ======================================================================
\section{Composite Channel Equations}
%% ======================================================================

Classical equivalent-roughness formulas designed for compound channels, applied here to patchy hillslopes.
Each formula computes a composite roughness $\nc$ from the spatially distributed Manning's~$n_i$ and (where needed) depth~$h_i$.

\medskip
\begin{table}[h!]
\centering
\renewcommand{\arraystretch}{1.6}
\begin{tabularx}{\textwidth}{@{}clX@{}}
\toprule
\# & Method & Formula \\
\midrule
1 & Lotter &
  $\displaystyle \nc = \frac{\bar{h}^{\,5/3}}%
  {\left\langle \dfrac{h_i^{\,5/3}}{n_i} \right\rangle}$ \\[4pt]

2 & Pavlovskii &
  $\displaystyle \nc = \sqrt{\avg{n_i^2}}$ \\

3 & Horton--Einstein &
  $\displaystyle \nc = \avg{n_i^{3/2}}^{\,2/3}$ \\

4 & Cox (arith.\ mean) &
  $\displaystyle \nc = \avg{n_i}$ \\

5 & Krishnamurthy--Chr. &
  $\displaystyle \nc = \exp\!\left(\frac{\avg{h_i^{3/2}\ln n_i}}%
  {\avg{h_i^{3/2}}}\right)$ \\[4pt]

6 & Colebatch &
  $\displaystyle \nc = \left(\frac{\avg{h_i\, n_i^{3/2}}}%
  {\bar{h}}\right)^{\!2/3}$ \\[4pt]

7 & Yen $R^{1/6}$ &
  $\displaystyle \nc = \bar{h}^{\,1/6}\left\langle
  \frac{n_i}{h_i^{\,1/6}} \right\rangle$ \\[4pt]

8 & Yen $R^{1/3}$ &
  $\displaystyle \nc = \frac{\avg{h_i^{\,1/3}\, n_i}}%
  {\bar{h}^{\,1/3}}$ \\[4pt]

9 & Felkel &
  $\displaystyle \nc = \frac{1}{\avg{n_i^{-1}}}$ \\
\bottomrule
\end{tabularx}
\caption{Composite-channel equations for equivalent roughness.
  $n_i$ = Manning's coefficient at cell~$i$;
  $h_i$ = depth at cell~$i$;
  $\bar{h} = \avg{h_i}$ = spatial mean depth;
  $\avg{\cdot}$ = spatial average.}
\label{tab:channel-eqs}
\end{table}

\subsection{Lotter Reduces to Unity Under CF Approximation}

Under the kinematic-wave correction-factor (CF) approximation, $h_i \propto n_i^{3/5}$, so
\[
  h_i^{5/3} = C \cdot n_i
  \quad\text{where } C = \bigl(q_0\, l / \So^{1/2}\bigr)
\]
is cell-independent.  Therefore $\avg{h^{5/3}/n} = C$ and $\nc = \bar{h}^{5/3}/(C\avg{n})$, which equals $\avg{n}$, making $\nc/\avg{n} = 1$ --- i.e., the Lotter equation becomes uninformative when depth variations follow the kinematic-wave scaling.


%% ======================================================================
\section{Length-Scale Attribution of $\Sf$ Decomposition}
%% ======================================================================

The six second-order terms $T_k$ from \eqref{eq:Sf-decomp} sum to a
departure score
\begin{equation}\label{eq:S-sum}
  S = \sum_{k=1}^{6} T_k
  = \avg{r^2} + \avg{\upsilon^2} + 4\avg{r\upsilon}
    - \tfrac{8}{3}\bigl(\avg{r\eta} + \avg{\upsilon\eta}\bigr)
    + \tfrac{14}{9}\avg{\eta^2}
\end{equation}
that controls how far the effect ratio departs from unity.
When $S$ varies across patch lengthscales ($L_V$), the following
methods identify which terms drive that variation.

\subsection{Covariance-Based Variance Attribution}

The key question is: \emph{when we vary the vegetation lengthscale
$L_V$ across simulations, which of the six terms $T_k$ are
responsible for the resulting spread in~$S$?}

Since $S = \sum_k T_k$, a standard identity from probability theory
gives an exact decomposition of variance:
\begin{equation}\label{eq:var-decomp}
  \operatorname{Var}(S)
  = \sum_{k=1}^{6} \operatorname{Cov}(T_k,\, S)
\end{equation}

Dividing both sides by $\operatorname{Var}(S)$ yields the
\emph{attribution fraction} for term~$k$:
\begin{equation}\label{eq:attrib-frac}
  \alpha_k
  = \frac{\operatorname{Cov}(T_k,\, S)}{\operatorname{Var}(S)},
  \qquad \sum_{k=1}^{6}\alpha_k = 1
\end{equation}

\paragraph{Interpretation of sign.}
\begin{itemize}
  \item $\alpha_k > 0$: term~$T_k$ co-varies with $S$ in the same
    direction --- it \emph{drives} the cross-scale variability.
    When $S$ is large, $T_k$ also tends to be large.
  \item $\alpha_k < 0$: term~$T_k$ moves opposite to $S$ ---
    it \emph{dampens} or partially offsets the other terms.
  \item The fractions always sum to exactly~1 by construction,
    regardless of sample size.
\end{itemize}

The cross-terms with negative coefficients in $S$ (namely
$-\frac{8}{3}\avg{r\eta}$ and $-\frac{8}{3}\avg{\upsilon\eta}$)
are the most likely candidates for negative $\alpha_k$, since their
weighted products naturally tend to oppose the sum.  The pure variance
terms ($\avg{r^2}$, $\avg{\eta^2}$) tend to stay positive because
variances are non-negative and generally increase with~$L_V$.

\paragraph{Grouping.}
The covariance in \eqref{eq:var-decomp} is computed across the
set of simulations that differ in~$L_V$ within a fixed group.
Two grouping strategies are used:
\begin{enumerate}
  \item \textbf{By anisotropy:} gradient-aligned, isotropic, and
    contour-aligned, with $f_V$ and storm forcing pooled within each group.
  \item \textbf{By vegetation fraction:} $f_V$ is binned into three
    equal-frequency classes (via quantile binning) to ensure sufficient
    sample sizes.  Anisotropy and storm forcing are pooled within
    each bin.
\end{enumerate}

\paragraph{Stability considerations.}
When the number of simulations per group is small ($n \lesssim 5$),
the covariance estimates can be noisy and individual $\alpha_k$ values
may exhibit spurious sign flips --- particularly for the cross-terms
whose true values are close to zero.  Pooling across all converged
storms (rather than conditioning on a single storm event) and binning
$f_V$ into terciles rather than using raw levels substantially
improves stability.

\subsection{Slope Decomposition (Fixed-Effect Regression)}

While the covariance attribution answers ``what fraction of spread
is due to each term?'', the slope decomposition asks a more direct
question: \emph{how much does each $T_k$ change per unit increase
in~$L_V$?}

Fit each term with an OLS model that includes $L_V$ as the
predictor of interest and dummy variables to absorb group-level
variation:
\begin{equation}\label{eq:slope-decomp}
  T_k = \beta_{0,k} + \beta_{L,k}\, L_V
        + \sum_j \gamma_{j,k}\,\mathbb{1}_{f_V=j}
        + \sum_m \delta_{m,k}\,\mathbb{1}_{\text{storm}=m}
        + \varepsilon_k
\end{equation}
where $\mathbb{1}_{f_V=j}$ are vegetation-fraction fixed effects,
$\mathbb{1}_{\text{storm}=m}$ are storm-forcing fixed effects
(encoding each $p \times t_r$ combination), and $\beta_{L,k}$ is
the marginal effect of lengthscale on term~$k$, holding all
else constant.

The slopes decompose the total sensitivity of $S$ to patch scale:
\begin{equation}\label{eq:slope-sum}
  \frac{\partial S}{\partial L_V}
  = \sum_{k=1}^{6} \beta_{L,k}
\end{equation}
Terms with large $|\beta_{L,k}|$ are the primary drivers of the
scale dependence of~$S$.  Standard errors from OLS provide
confidence intervals that are unavailable from the model-free
covariance attribution.

When the regression is run separately per group (e.g.\ per
anisotropy level or per $f_V$ bin), the fixed effects absorb only
the remaining nuisance variables --- for instance, when grouping by
anisotropy, the design matrix includes $f_V$ and storm fixed effects
but not anisotropy dummies.

\subsection{Summary: Two Complementary Views}

\begin{center}
\begin{tabular}{@{}lll@{}}
\toprule
Analysis & Question & Method \\
\midrule
Covariance attribution
  & Which terms explain the \emph{spread} in $S$?
  & Variance identity (exact, model-free) \\
Slope decomposition
  & How much does each term change per metre of $L_V$?
  & OLS regression (gives standard errors) \\
\bottomrule
\end{tabular}
\end{center}


\end{document}
